% Автоматаар үүсгэсэн — эх файл: docs/report-template.md
\chapter{Тайлангийн загвар}

\textbf{Баг:} ule{1.5cm}{0.4pt} \textbf{Огноо:} ule{1.5cm}{0.4pt}
\textbf{Гишүүд ба хувь нэмэр:}

\begingroup\scriptsize\begin{longtable}[]{@{}ll@{}}
\toprule
Гишүүн & Хийсэн ажил \\
\midrule
\endhead
\bottomrule
\endlastfoot
& \\
& \\
& \\
\end{longtable}\endgroup

\textbf{Git commit:} `ule{1.5cm}{0.4pt}\texttt{**Tag:**}lab0\_\_-done`

\hypertarget{ux445ux44dux43cux436ux438ux43bux442ux438ux439ux43d-ux43eux440ux447ux438ux43d}{%
\section{0. Хэмжилтийн орчин}\label{ux445ux44dux43cux436ux438ux43bux442ux438ux439ux43d-ux43eux440ux447ux438ux43d}}

Тоон үр дүн зөвхөн орчноо мэдэж байж давтагдана. Хэмжилт бүрийн өмнө бөглөнө.

\begingroup\scriptsize\begin{longtable}[]{@{}lll@{}}
\toprule
& {[}Pi{]} Raspberry Pi 3B & {[}ЗК{]} Зөөврийн компьютер \\
\midrule
\endhead
\bottomrule
\endlastfoot
Загвар / процессор & & \\
RAM (нийт / боломжтой) & & \\
OS, цөмийн хувилбар & & \\
Сүлжээ (кабель / Wi-Fi, хурд) & & \\
Хадгалалт (microSD ангилал) & & \\
Docker хувилбар & & \\
\end{longtable}\endgroup

\textbf{Хэмжилтийн өмнөх шалгалт ({[}Pi{]}):}

\begingroup\scriptsize\begin{longtable}[]{@{}lll@{}}
\toprule
Шалгалт & Утга & Хэвийн үү \\
\midrule
\endhead
\bottomrule
\endlastfoot
\texttt{vcgencmd\ get\_\allowbreak{}throttled} & & \texttt{0x0} байх ЁСТОЙ \\
\texttt{vcgencmd\ measure\_\allowbreak{}temp} & & \textless{} 70 °C \\
\texttt{free\ -m} → available & & \textgreater{} 400 MiB \\
\texttt{timedatectl} synchronized & & yes \\
VS Code сервер таслагдсан уу & & \texttt{pkill\ -f\ vscode-server} \\
\end{longtable}\endgroup

\begin{quote}
⚠ \texttt{get\_throttled} нь \texttt{0x0} биш байсан бол тэр хэмжилт \textbf{хүчингүй}. Хөргөөд дахин хий, тайланд хуучин тоог бичиж болохгүй.
\end{quote}

\hypertarget{ux433ux4afux439ux446ux44dux442ux433ux44dux441ux44dux43d-ux430ux436ux438ux43b}{%
\section{1. Гүйцэтгэсэн ажил}\label{ux433ux4afux439ux446ux44dux442ux433ux44dux441ux44dux43d-ux430ux436ux438ux43b}}

Товч, алхмаар. Дэлгэцийн зураг биш --- \textbf{комманд ба гаралт}.

\begin{verbatim}
# жишээ
$ docker compose --profile core up -d
...
\end{verbatim}

\hypertarget{ux445ux44dux43cux436ux438ux43bux442ux438ux439ux43d-ux4afux440-ux434ux4afux43d}{%
\section{2. Хэмжилтийн үр дүн}\label{ux445ux44dux43cux436ux438ux43bux442ux438ux439ux43d-ux4afux440-ux434ux4afux43d}}

\begin{quote}
Зааврын хүснэгтүүдийг энд хуулж бөглөнө. \textbf{Бөглөөгүй хүснэгт = оноо хасагдана.}
\end{quote}

\hypertarget{ux434ux4afux433ux43dux44dux43bux442}{%
\section{3. Дүгнэлт}\label{ux434ux4afux433ux43dux44dux43bux442}}

Тоон үр дүн юу гэсэн үг вэ? Багийн төслийн шийдэлд хэрхэн нөлөөлөх вэ?
Хэмжилтэд суурилаагүй дүгнэлт оноо авахгүй.

\hypertarget{ux445ux44fux43dux430ux43bux442ux44bux43d-ux430ux441ux443ux443ux43bux442ux44bux43d-ux445ux430ux440ux438ux443ux43bux442}{%
\section{4. Хяналтын асуултын хариулт}\label{ux445ux44fux43dux430ux43bux442ux44bux43d-ux430ux441ux443ux443ux43bux442ux44bux43d-ux445ux430ux440ux438ux443ux43bux442}}

\begin{enumerate}
\def\labelenumi{\arabic{enumi}.}
\tightlist
\item
\item
\item
\end{enumerate}

\hypertarget{ux442ux443ux43bux433ux430ux440ux441ux430ux43d-ux430ux441ux443ux443ux434ux430ux43b-ux431ux430-ux448ux438ux439ux434ux44dux43b}{%
\section{5. Тулгарсан асуудал ба шийдэл}\label{ux442ux443ux43bux433ux430ux440ux441ux430ux43d-ux430ux441ux443ux443ux434ux430ux43b-ux431ux430-ux448ux438ux439ux434ux44dux43b}}

\begingroup\scriptsize\begin{longtable}[]{@{}llll@{}}
\toprule
Асуудал & Шалтгаан & Шийдсэн арга & Хугацаа \\
\midrule
\endhead
\bottomrule
\endlastfoot
& & & \\
\end{longtable}\endgroup

\hypertarget{ux430ux448ux438ux433ux43bux430ux441ux430ux43d-ux44dux445-ux441ux443ux440ux432ux430ux43bux436}{%
\section{6. Ашигласан эх сурвалж}\label{ux430ux448ux438ux433ux43bux430ux441ux430ux43d-ux44dux445-ux441ux443ux440ux432ux430ux43bux436}}

\begin{itemize}
\tightlist
\item
  Баримт бичиг, өгүүлэл, форум
\item
  Хиймэл оюуны туслах хэрэгсэл ашигласан бол: хэзээ, юунд, хэрхэн шалгасан
\end{itemize}

\hypertarget{ux445ux44dux43cux436ux438ux43bux442ux438ux439ux43d-ux444ux430ux439ux43bux443ux443ux434}{%
\section{7. Хэмжилтийн файлууд}\label{ux445ux44dux43cux436ux438ux43bux442ux438ux439ux43d-ux444ux430ux439ux43bux443ux443ux434}}

\texttt{measurements/} эсвэл \texttt{labNN/out/} доторх файлын нэрс (файлыг Git-д оруулахгүй):

\begin{itemize}
\tightlist
\item
\end{itemize}
